\emph{Setting.} The check uses the five-mode module of Proposition \claimref{standard-model.majorana-element} as $32 \times 32$ rational matrices (creation $v_{m}$, annihilation $w_{m}$, Jordan-Wigner signs), with the real-form Majoranas of that proposition's check, $\gamma_{2m} = v_{m} + w_{m}$ and $\gamma_{2m+1} = w_{m} - v_{m}$, which anticommute and square to $+1$ and $-1$: the module is the irreducible representation of $\Cl(5,5) \cong M_{32}(\R)$, so the scalar part of a Clifford element is the normalised trace, $\langle A\rangle_{0} = \tfrac{1}{32}\mathrm{tr}\,A$, and the Clifford bilinear form is $Q(A, B) = \langle A\tilde{B}\rangle_{0}$. On a five-blade reversion is the identity, since reversing five anticommuting factors is $\binom{5}{2} = 10$ transpositions (the check verifies $\gamma_{p_{5}}\cdots\gamma_{p_{1}} = \gamma_{p_{1}}\cdots\gamma_{p_{5}}$ on four blades), so on grade five $Q(A, B) = \langle AB\rangle_{0}$. This is the form whose value on the frozen point is the frozen-point condition: Fact \claimref{mass.imported-frozen-point-tractor-connection} gives $Q(I) = \Lambda/3$ for the parallel tractor $I = \nInf + \tfrac{\Lambda}{6}\nO$ of a vacuum Einstein field, a real vector of the real representation $\R^{2,4}$, and the stabiliser and the de Sitter radius are read off that norm. The question of the open problem's third line is whether an invariant condition of this kind, $Q(T) = M^{2}$ or any invariant algebraic equation on the element alone, can fix the vev along the Majorana element.

\emph{(i)} $T = v_{1}v_{2}v_{3}v_{4}v_{5}$ with $v_{m}^{2} = 0$ and $v_{m}v_{n} = -v_{n}v_{m}$, so $T^{2} = \pm v_{1}v_{1}\cdots = 0$: $T$ is a zero divisor and $Q(T) = \langle T^{2}\rangle_{0} = 0$. The check verifies $T^{2} = 0$, $Q(T) = 0$, and as a control that the non-null blade $\gamma_{0}\gamma_{2}\gamma_{4}\gamma_{6}\gamma_{8}$ has $Q = 1$. The Fock adjoint is the matrix transpose on this rational module, $T^{\dagger} = w_{5}\cdots w_{1}$, and $TT^{\dagger}$ is the projector on the full state, so $\langle TT^{\dagger}\rangle_{0} = \tfrac{1}{32}$ (checked entrywise). $[B - L, T] = -2T$ is clause (2) of Proposition \claimref{standard-model.majorana-element}, re-verified with $B - L = 1 - \tfrac23 n_{W}$.

\emph{(ii)} In $\Cl(p,q)$ with $n = p + q = 10$ the pseudoscalar satisfies $I_{10}^{2} = (-1)^{n(n-1)/2}(-1)^{q} = (-1)^{45}(-1)^{5} = +1$ and anticommutes with every vector since $n$ is even (both checked). For $A, B$ of odd grade, $AI_{10} = -I_{10}A$, so $I_{10}A\,I_{10}B = -I_{10}^{2}AB = -AB$ and $Q(I_{10}A, I_{10}B) = -Q(A, B)$: an anti-isometry. If $I_{10}A = \epsilon A$ and $I_{10}B = \epsilon B$ with $\epsilon = \pm1$ then $Q(A, B) = Q(I_{10}A, I_{10}B)/\epsilon^{2}\cdot(-1)$, that is $Q(A, B) = -Q(A, B) = 0$. The two eigenspaces of $I_{10}$ on the $252$ five-vectors have dimension $126$ each (the check computes the rank of $\{A + I_{10}A\}$ over the $252$ basis blades and finds $126$); over $\C$ they are the $126$ and $\overline{126}$ of $\mathfrak{so}(10)$, the irreducible representations of self-dual and anti-self-dual five-forms (standard, cited from memory), and the $(10,1,3)$ orbit of $T$ in Proposition \claimref{standard-model.majorana-element}(3) lies in the one containing $T$. The check verifies $I_{10}T = T$, the anti-isometry on three sampled five-vectors of nonzero norm, $Q(A_{\pm}) = 0$ for their projections $A_{\pm} = \tfrac12(A \pm I_{10}A)$, and $Q(A_{+}, A_{-}) = \tfrac12 Q(A)$, so $Q$ is nondegenerate on grade five and pairs the two halves. $T^{\dagger}$ lies in the $-1$ eigenspace and $Q(T, T^{\dagger}) = \tfrac{1}{32}$: the Fock norm of $T$ is the Clifford pairing of $T$ with an element of the other half.

\emph{(iii)} Let $P$ be a polynomial function on the five-vectors invariant under a group containing the one-parameter subgroup $\exp(i\theta(B - L))$ of the compact Pati-Salam group (the compact real form is spanned by $i$ times the Hermitian elements of the $21$, as in the stored check; $B - L$ is Hermitian). Since $[B - L, T] = -2T$, $\exp(i\theta(B - L))\,T\,\exp(-i\theta(B - L)) = e^{-2i\theta}T$, so $P(T) = P(e^{-2i\theta}T) = \sum_{d}e^{-2i\theta d}P_{d}(T)$ for all $\theta$, where $P_{d}$ is the homogeneous component of degree $d$; hence $P_{d}(T) = 0$ for every $d \ge 1$ and $P(T) = P(0)$. Invariance gives $P(gTg^{-1}) = P(T)$ on the whole orbit. The same holds for the real form: $B - L = 1 - \tfrac23 n_{W}$ is a real diagonal matrix, an element of the real $\mathfrak{so}(5,5)$ represented on this module (the $v_{m}$, $w_{m}$ are real), so $\exp(t(B - L))$ for real $t$ lies in the real spin group and $\exp(t(B - L))\,T\,\exp(-t(B - L)) = e^{-2t}T$; a continuous invariant $P$ of the real form then has $P(T) = P(e^{-2t}T) \to P(0)$ as $t \to \infty$: $T$ lies in the null cone (the Hilbert-Mumford picture, standard). This says nothing about the whole $126$: its invariant ring is not trivial (the group has dimension $45$ and the representation $126$, so no orbit is dense and nonconstant invariants exist; standard, cited from memory), and they vanish on $T$ and its orbit, not on the whole $126$. Invariant bilinear forms are a separate matter: the $126$ of $\mathfrak{so}(10)$ is not isomorphic to its dual (its weights are not closed under negation, as $T$ has weight $(\tfrac12, \ldots, \tfrac12)$ in the Cartan of the $n_{m} - \tfrac12$ while $T^{\dagger}$, of the opposite weight, lies in the $-1$ eigenspace of $I_{10}$, checked), so it carries no invariant bilinear form at all, Clifford or other; the vanishing of $Q$ on it in (ii) is the instance of this that the series' form exhibits. The check runs this for $Q$ on five rotors $g$ of the $21$ ($1 + X$ with $X$ nilpotent: the $SU(2)_{R}$ lowering operator $w_{5}w_{4}$, a pair creator $v_{1}v_{2}$, the $SU(2)_{L}$ raising operator $v_{4}w_{5}$, a pair annihilator $w_{1}w_{3}$, and a product of two), three of which move $T$: $gTg^{-1}$ stays in the $+1$ eigenspace, $Q(gTg^{-1}) = 0$ and $Q(gTg^{-1}, T) = 0$.

\emph{(iv)} The adjoint fixes $\gamma_{2m}$ and negates $\gamma_{2m+1}$ (checked for all $m$): it is reversion composed with the reflection of the five negative-square generators. That reflection is not in the orthogonal group of the ten, so the adjoint is not an anti-automorphism of $\Cl(5,5)$ as a Clifford algebra; it is determined by the splitting of the ten into the span of the $\gamma_{2m}$ and the span of the $\gamma_{2m+1}$, equivalently by the maximal isotropic subspace spanned by the $v_{m}$, the particle modes, and this is the Hermitian structure of the module (the structure through which, in the register's rung-unit premise synthesis.rung-unit-second-free-number, cited by name and not assumed, the rung unit is seen while the bilinear tractor norm does not see it). A sesquilinear condition $\langle TT^{\dagger}\rangle_{0}|\tau|^{2} = M^{2}$ on a vev $\tau T$ fixes $|\tau|$ against $M$ and leaves $M$ as the input. Proposition \claimref{synthesis.majorana-scale-free-on-pati-salam-route} shows the stored relations fix that input only through the observed $m_{3}$; this result shows the frozen-point mechanism, an invariant norm of an invariant element, has no counterpart that could fix it either. The contrast with $\Lambda$ is exact: the frozen point is a real vector with $Q(I) = \Lambda/3 \neq 0$, and its norm is the scale; the Majorana element has $Q(T) = 0$ identically, on its orbit and on its whole irreducible, and no holomorphic invariant distinguishes a vev along it from zero.
